\section{逆解析、測定誤差および識別可能性}
\subsection{推定する状態量と観測の共有}
推定対象$\boldsymbol\theta$は、温度、電子密度、選択した外部粒子密度など、設定ファイルの明示的な経路へ結びつける。正の密度等には$x=\log_{10}\theta$を使え、上下限もその座標へ変換する。報告値は名前付きの物理単位へ戻す。

複数の装置、視線、繰返し測定を同じプラズマ状態へ接続できる。観測IDと予測IDを分離するため、三回の観測を一つの予測スペクトルへ対応させることも可能である。現行の一回のfitは、一つの\code{condition\_id}と一つの\code{state\_group\_id}を扱う。異なる電力・圧力条件を同時に推定する階層モデルとは区別する。

\subsection{残差と共分散}
選択した画素の観測を$\boldsymbol y$、予測を$\boldsymbol F(\boldsymbol\theta)$とする。測定共分散$C=LL^\mathsf T$に対し、
\begin{equation}
\boldsymbol r_{\mathrm{data}}(\boldsymbol\theta)
=L^{-1}\left[\boldsymbol F(\boldsymbol\theta)-\boldsymbol y\right]
\label{eq:whiten}
\end{equation}
と白色化する。独立誤差なら$r_i=(F_i-y_i)/\sigma_i$に戻る。共分散は有限・対称・正定値を要し、明示的な逆行列ではなくCholesky因子で解く。測定$\sigma$も共分散もない相対fitでは、観測の標準偏差等による便宜的尺度を使う。この場合のcostを、校正済みの統計的$\chi^2$と同じ意味で扱うことはできない。

同じ校正を共有する観測には、共通の乗算誤差がある。校正group $g$に属する観測$i$の相対標準不確かさを$s_i$、観測値を使った固定ベクトルを$\boldsymbol u_g$として
\begin{equation}
C=C_{\mathrm{meas}}+\sum_g\boldsymbol u_g\boldsymbol u_g^\mathsf T,
\qquad (u_g)_i=\begin{cases}s_iy_i,&i\in g,\\0,&i\notin g\end{cases}
\label{eq:calcov}
\end{equation}
とする。予測値で共分散を毎回更新するのではなく、観測を尺度基準として固定する。共通誤差を独立画素誤差とみなしてしまうと、画素数や繰返し数に比例して精度が向上したように見えるため、この区別が重要になる。

測定の画素状態を保持し、通常残差には有効な画素を使う。飽和画素には下限値のメタデータを要求するが、現行の目的関数では有効画素以外を除外し、飽和の打切り尤度や一側制約を評価する機能は実装していない。入力adapterも、校正済みCSVの単位、誤差、状態、出典を検査して接続するもので、raw countsの校正や平滑化を暗黙に行わない。

\subsection{目的関数、波長選択、補助パラメータ}
推定は、測定由来の残差、事前分布、殻方向平滑化を組み合わせた
\begin{equation}
\widehat{\boldsymbol\theta}=\mathop{\arg\min}_{\boldsymbol\theta\in\Theta}
\frac12\left(\|\boldsymbol r_{\mathrm{data}}\|_2^2+
\|\boldsymbol r_{\mathrm{prior}}\|_2^2+
\|\boldsymbol r_{\mathrm{smooth}}\|_2^2\right)
\label{eq:objective}
\end{equation}
で行う。全スペクトル、選択窓内の形、baseline補正後の面積・peak・比、名前付きfeatureの共分散を使える。面積や比の残差は所定の尺度で無次元化してから白色化する。同じfeatureを共分散項とスカラー重み項に二重に入れない。ただし、同じ画素から作ったfeatureと全画素残差を別々に足した目的関数は、自動的に独立な尤度にはならない。採用した重みの意味と情報の重複を吟味する。

通常の\code{full\_spectrum}は全有効画素を使う。\code{window\_union}を選ぶと、物理的に対応づけたfit窓の和集合だけを使える。\code{score}窓は推定後の予測確認へ保留し、fit窓との区間重複を禁止する。弱い窓を閾値で除く機能はあるが、共分散で名前を指定したfeatureが欠落した場合には、残差の長さを黙って変えずエラーとする。

相対fitでは、gain、offset、必要なら波長方向のtiltを内側で推定できる。これらも最終残差と同じ有効画素と白色化を使う。スペクトルのみなら重み付き線形最小二乗、featureや事前項を含む場合は対応する内側評価を行う。固定装置gain・baselineは順解析で一度だけ適用する。外側の最適化は、必要に応じたdifferential evolutionの後、境界付き\code{least\_squares}で局所精密化する。

\subsection{四つの推定モードとactinometry}
\begin{table}[htbp]
\centering\small
\begin{tabularx}{\linewidth}{p{35mm}YY}
\toprule
モード & 主な情報 & 定量解釈に必要な条件 \\
\midrule
\code{relative\_shape} & 相対形状、制約された空間変化 & gainや供給密度との尺度縮退を確認する \\
\code{ratio\_diagnostic} & 有効な窓面積・peak比 & 波長間応答と比の動力学を与える \\
\code{actinometry} & 同じ比目的関数 & actinometerと対象の励起・消光・組成の前提を別途支持する \\
\code{calibrated\_absolute} & 絶対尺度の測定 & 物理発光、校正基準の一致、固定gain等、測定誤差、供給密度と損失の独立制約 \\
\bottomrule
\end{tabularx}
\caption{モードは利用目的と解釈の契約を明示する。選択した名前だけで、温度や密度の観測可能性が得られるわけではない。}
\end{table}

actinometryは、既知の参照粒子Aと対象Xの発光比から密度比を推定する考え方である\cite{coburn,fuller}。視線内で一様な状態、同じ電子密度・光路、光学的薄さ、直接励起、無視できる消光を仮定した単純なコロナ極限では
\begin{equation}
\frac{I_X}{I_A}=\frac{N_X}{N_A}\frac{k_X}{k_A}
\frac{B_X}{B_A}\frac{\lambda_A}{\lambda_X}\frac{T(\lambda_X)}{T(\lambda_A)}
\label{eq:actino}
\end{equation}
と書ける。$B$は放射分岐割合である。波長応答を補正した強度なら最後の透過率比は除く。実際には段階励起、解離励起、消光、EEDF変化の採否を確認しなければならない。OESCRの\code{actinometry}は、この前提を明示して同じ比の最適化を利用する。別の化学組成solverを起動するモードではない。

絶対モードでは自動gain・offset・tiltや装置尺度の同時推定を禁止し、測定CSVと装置の基準・単位・校正参照を照合する。しかし、$I\propto\nelec N_s k(\Te)$の$\nelec$と$N_s$をともに自由にすれば、校正済みでも積しか分からない場合がある。絶対校正の必要条件と、パラメータの識別条件は分けて判定する。

\subsection{SVD、profile、局所不確かさの読み方}
局所識別性は、事前分布や平滑化を除いた測定残差だけから、最適化座標$\boldsymbol x$に対する有限差分Jacobian
\begin{equation}
J_{ij}=\frac{\partial r_{\mathrm{data},i}}{\partial x_j},\qquad
J=U\,\operatorname{diag}(s_1,\ldots,s_p)V^\mathsf T
\label{eq:svd}
\end{equation}
を作って調べる。小さい特異値に対応する$V$の方向は、観測をほとんど変えずに動かせるパラメータの組合せである。条件数$s_{\max}/s_{\min}$とrankは座標尺度・局所点・閾値に依存する。full rankでも実用上の精度を保証しない。

明示したnuisanceパラメータのJacobianを$J_\eta$、targetを$J_\theta$とすると、
\begin{equation}
J_{\theta,\perp}=(I-Q_\eta Q_\eta^\mathsf T)J_\theta
\label{eq:project}
\end{equation}
で、nuisanceが説明できる測定方向を除く。$Q_\eta$はSVDによる列空間の正規直交基底である。内側で既にprofileしたgain等は二度数えない。事前分布を入れてrankを増やしたように見せることも避ける。

非線形性は、対象パラメータ$\psi$を固定して残り$\boldsymbol\lambda$を再推定する
\begin{equation}
D_p(\psi)=\min_{\boldsymbol\lambda}\|\boldsymbol r_{\mathrm{data}}(\psi,\boldsymbol\lambda)\|_2^2
\label{eq:profile}
\end{equation}
でも調べられる。局所的で境界付きのprofileであり、score窓、事前項、平滑化を除く。狭い閉じた谷か、広い谷か、境界へ開いているかを確認する。$\Delta D=1$や3.84は追加条件下の漸近的目安であり、現行の図から校正済み信頼区間と呼ぶことはしない。

Laplace近似は全目的関数の局所曲率から共分散を近似する。事前・平滑化・指定共分散に条件づけられ、断面積、欠落反応、モデル不一致の不確かさを周辺化した完全な事後分布ではない。これらの監査は\code{--audit}等で明示的に実行し、毎回の順解析反復には付けない。最適化の\code{success}、数値品質、実用的識別性、外部検証は、それぞれ別の判定である。
